\documentclass[10pt,a4paper]{amsart}
\usepackage{amsmath,amssymb,amsthm,graphicx}
\usepackage{verbatim}
\newenvironment{reference}{\comment}{\endcomment}
%\newenvironment{reference}{}{}
\newenvironment{slogan}{\comment}{\endcomment}
\newenvironment{history}{\comment}{\endcomment}

% For commutative diagrams we use Xy-pic
\usepackage[all]{xy}

% We use 2cell for 2-commutative diagrams.
\xyoption{2cell}
% Theorem environments.
%
\theoremstyle{plain}
\newtheorem{theorem}[subsection]{Theorem}
\newtheorem{proposition}[subsection]{Proposition}
\newtheorem{lemma}[subsection]{Lemma}

\theoremstyle{definition}
\newtheorem{definition}[subsection]{Definition}
\newtheorem{example}[subsection]{Example}
\newtheorem{exercise}[subsection]{Exercise}
\newtheorem{situation}[subsection]{Situation}

\theoremstyle{remark}
\newtheorem{remark}[subsection]{Remark}
\newtheorem{remarks}[subsection]{Remarks}

\numberwithin{equation}{subsection}

% Macros
%
\def\lim{\mathop{\mathrm{lim}}\nolimits}
\def\colim{\mathop{\mathrm{colim}}\nolimits}
\def\Spec{\mathop{\mathrm{Spec}}}
\def\Hom{\mathop{\mathrm{Hom}}\nolimits}
\def\Ext{\mathop{\mathrm{Ext}}\nolimits}
\def\SheafHom{\mathop{\mathcal{H}\!\mathit{om}}\nolimits}
\def\SheafExt{\mathop{\mathcal{E}\!\mathit{xt}}\nolimits}
\def\Sch{\mathit{Sch}}
\def\Mor{\mathop{\mathrm{Mor}}\nolimits}
\def\Ob{\mathop{\mathrm{Ob}}\nolimits}
\def\Sh{\mathop{\mathit{Sh}}\nolimits}
\def\NL{\mathop{N\!L}\nolimits}
\def\CH{\mathop{\mathrm{CH}}\nolimits}
\def\proetale{{pro\text{-}\acute{e}tale}}
\def\etale{{\acute{e}tale}}
\def\QCoh{\mathit{QCoh}}
\def\Ker{\mathop{\mathrm{Ker}}}
\def\Im{\mathop{\mathrm{Im}}}
\def\Coker{\mathop{\mathrm{Coker}}}
\def\Coim{\mathop{\mathrm{Coim}}}

% Boxtimes
%
\DeclareMathSymbol{\boxtimes}{\mathbin}{AMSa}{"02}

%
% Macros for moduli stacks/spaces
%
\def\QCohstack{\mathcal{QC}\!\mathit{oh}}
\def\Cohstack{\mathcal{C}\!\mathit{oh}}
\def\Spacesstack{\mathcal{S}\!\mathit{paces}}
\def\Quotfunctor{\mathrm{Quot}}
\def\Hilbfunctor{\mathrm{Hilb}}
\def\Curvesstack{\mathcal{C}\!\mathit{urves}}
\def\Polarizedstack{\mathcal{P}\!\mathit{olarized}}
\def\Complexesstack{\mathcal{C}\!\mathit{omplexes}}
% \Pic is the operator that assigns to X its picard group, usage \Pic(X)
% \Picardstack_{X/B} denotes the Picard stack of X over B
% \Picardfunctor_{X/B} denotes the Picard functor of X over B
\def\Pic{\mathop{\mathrm{Pic}}\nolimits}
\def\Picardstack{\mathcal{P}\!\mathit{ic}}
\def\Picardfunctor{\mathrm{Pic}}
\def\Deformationcategory{\mathcal{D}\!\mathit{ef}}
\usepackage{cleveref}
\crefname{originallemma}{original lemma}{original lemmas}
\Crefname{originallemma}{Original lemma}{Original lemmas}
\makeatletter
\newlabel{lemma-Japanese-local}{{032H}{source}{}{sourcepointer-1}{}}
\newlabel{lemma-Japanese-local@cref}{{[originallemma][0][]032H}{[1][0][]source}{}{}{}}
\makeatother
\crefname{originallemma}{original lemma}{original lemmas}
\Crefname{originallemma}{Original lemma}{Original lemmas}
\makeatletter
\newlabel{lemma-integral-closure-transitive}{{0308}{source}{}{sourcepointer-2}{}}
\newlabel{lemma-integral-closure-transitive@cref}{{[originallemma][0][]0308}{[1][0][]source}{}{}{}}
\makeatother
\crefname{originallemma}{original lemma}{original lemmas}
\Crefname{originallemma}{Original lemma}{Original lemmas}
\makeatletter
\newlabel{lemma-integral-permanence}{{02JM}{source}{}{sourcepointer-3}{}}
\newlabel{lemma-integral-permanence@cref}{{[originallemma][0][]02JM}{[1][0][]source}{}{}{}}
\makeatother
\crefname{originallemma}{original lemma}{original lemmas}
\Crefname{originallemma}{Original lemma}{Original lemmas}
\makeatletter
\newlabel{lemma-isolated-point-fibre}{{00PK}{source}{}{sourcepointer-4}{}}
\newlabel{lemma-isolated-point-fibre@cref}{{[originallemma][0][]00PK}{[1][0][]source}{}{}{}}
\makeatother
\crefname{originallemma}{original lemma}{original lemmas}
\Crefname{originallemma}{Original lemma}{Original lemmas}
\makeatletter
\newlabel{lemma-localize-N}{{032G}{source}{}{sourcepointer-5}{}}
\newlabel{lemma-localize-N@cref}{{[originallemma][0][]032G}{[1][0][]source}{}{}{}}
\makeatother
\crefname{originallemma}{original lemma}{original lemmas}
\Crefname{originallemma}{Original lemma}{Original lemmas}
\makeatletter
\newlabel{lemma-quasi-finite-open-integral-closure}{{00QB}{source}{}{sourcepointer-6}{}}
\newlabel{lemma-quasi-finite-open-integral-closure@cref}{{[originallemma][0][]00QB}{[1][0][]source}{}{}{}}
\makeatother
\title{The N-2 property ascends along a quasi-finite extension of Noetherian domains}
\author{The Stacks Project Authors}
\date{stacks-project:algebra.tex:local-snapshot-20261012; GNU FDL 1.2 or later}
\begin{document}
\maketitle
\noindent\textit{Editorial note.} All delivered mathematical text is the author's own, copied from the pinned Stacks Project chapter file algebra.tex; no statement or proof step was rewritten, reordered or removed.\par\medskip
\noindent\textit{Editorial note.} Preamble provenance, as required for this source: the frozen snapshot stages the chapter's own preamble file beside it (preamble.tex, 217 lines, 8116 bytes, the same shared Stacks-wide preamble the published chow and etale-cohomology cards use); the first line of the pinned chapter file is \textbackslash{}input\{preamble\}, so no preamble text lives inside the chapter itself. The card takes that file's lines 8-18 (the verbatim package and the three \textbackslash{}comment-based environments reference/slogan/history at lines 9-12, section 13.2, together with the Xy-pic setup) and its lines 154 to the runtime end of the file (the native theorem declarations and mathematical macros). Not one character was changed, no package was added, no macro was defined, and the excerpt serves only as preamble, never as a span coordinate.\par\medskip
\noindent\textit{Editorial note.} The same-chapter pointers are located at their \textbackslash{}label lines inside the pinned chapter file; the cross-chapter pointers are located in the project's own tag list staged in the snapshot (sources/stacks-chow/tags.txt), at the verbatim lines carrying their canonical tags.\par\medskip
\noindent\textit{Editorial note.} The counted proof refers to several same-chapter statements by label; they are recorded as navigation stubs pointed at their \textbackslash{}label lines, and the arguments consuming them stay inside the counted span.\par\medskip
\noindent\textit{Editorial note.} The counted spans contain no \textbackslash{}cite, so the card carries no bibliography.\par\medskip
\noindent\textit{Allowed standard prerequisites (editorial).}\begin{itemize}
\item Quasi-finiteness gives finiteness of the induced fraction field extension. Used at the start of the counted proof (lines 45425-45429). Provides the finite field extension that the N-2 property of the base is applied to.
\item The integral closure as a finite module over an N-2 base. Used in the middle of the counted proof (lines 45430-45446). Produces a finite intermediate ring whose integral closure over the target is the object whose finiteness has to be shown.
\end{itemize}\medskip
\ifdefined\hypertarget\hypertarget{sourcepointer-1}{}\fi
\noindent\textit{Source reference.} Original lemma 032H: sources/stacks-algebra/pinned-original-source.tex, source lines 45402--45402. Navigation pointer for \textbackslash{}ref\{lemma-Japanese-local\} used inside the counted span (line 45438); the argument that consumes it is retained in the span.\par\medskip
\ifdefined\hypertarget\hypertarget{sourcepointer-2}{}\fi
\noindent\textit{Source reference.} Original lemma 0308: sources/stacks-algebra/pinned-original-source.tex, source lines 7940--7940. Navigation pointer for \textbackslash{}ref\{lemma-integral-closure-transitive\} used inside the counted span (line 45446); the argument that consumes it is retained in the span.\par\medskip
\ifdefined\hypertarget\hypertarget{sourcepointer-3}{}\fi
\noindent\textit{Source reference.} Original lemma 02JM: sources/stacks-algebra/pinned-original-source.tex, source lines 7927--7927. Navigation pointer for \textbackslash{}ref\{lemma-integral-permanence\} used inside the counted span (line 45448); the argument that consumes it is retained in the span.\par\medskip
\ifdefined\hypertarget\hypertarget{sourcepointer-4}{}\fi
\noindent\textit{Source reference.} Original lemma 00PK: sources/stacks-algebra/pinned-original-source.tex, source lines 30413--30413. Navigation pointer for \textbackslash{}ref\{lemma-isolated-point-fibre\} used inside the counted span (line 45427); the argument that consumes it is retained in the span.\par\medskip
\ifdefined\hypertarget\hypertarget{sourcepointer-5}{}\fi
\noindent\textit{Source reference.} Original lemma 032G: sources/stacks-algebra/pinned-original-source.tex, source lines 45390--45390. Navigation pointer for \textbackslash{}ref\{lemma-localize-N\} used inside the counted span (line 45438); the argument that consumes it is retained in the span.\par\medskip
\ifdefined\hypertarget\hypertarget{sourcepointer-6}{}\fi
\noindent\textit{Source reference.} Original lemma 00QB: sources/stacks-algebra/pinned-original-source.tex, source lines 31123--31123. Navigation pointer for \textbackslash{}ref\{lemma-quasi-finite-open-integral-closure\} used inside the counted span (line 45434); the argument that consumes it is retained in the span.\par\medskip
\begin{lemma}
\label{lemma-quasi-finite-over-Noetherian-japanese}
Let $R$ be a domain. Let $R \subset S$ be a quasi-finite extension of domains
(for example finite). Assume $R$ is N-2 and Noetherian. Then $S$ is N-2.
\end{lemma}

\begin{proof}
Let $L/K$ be the induced extension of fraction fields.
Note that this is a finite field extension (for example by
Lemma \ref{lemma-isolated-point-fibre} (2)
applied to the fibre $S \otimes_R K$, and the definition of a
quasi-finite ring map).
Let $S'$ be the integral closure of $R$ in $S$.
Then $S'$ is contained in the integral closure of $R$ in $L$
which is finite over $R$ by assumption. As $R$ is Noetherian this
implies $S'$ is finite over $R$.
By Lemma \ref{lemma-quasi-finite-open-integral-closure}
there exist elements $g_1, \ldots, g_n \in S'$
such that $S'_{g_i} \cong S_{g_i}$ and such that $g_1, \ldots, g_n$
generate the unit ideal in $S$. Hence it suffices to show that
$S'$ is N-2 by Lemmas \ref{lemma-localize-N} and \ref{lemma-Japanese-local}.
Thus we have reduced to the case where $S$ is finite over $R$.

\medskip\noindent
Assume $R \subset S$ with hypotheses as in the lemma and moreover
that $S$ is finite over $R$. Let $M$ be a finite field extension
of the fraction field of $S$. Then $M$ is also a finite field extension
of $K$ and we conclude that the integral closure $T$ of $R$ in
$M$ is finite over $R$. By Lemma \ref{lemma-integral-closure-transitive}
we see that $T$ is also the integral closure of $S$ in $M$ and we win by
Lemma \ref{lemma-integral-permanence}.
\end{proof}

\end{document}
